Note that we have only to prove (2) for $n\ge4$ and (3) for $n\ge5$. 

We first remark that in our cases the similarity matrix $Z_X(t)$ is invertible. 
When $n=3$ the similarity matrix $Z_X(t)$ is invertible for any $t$ $(t>0)$ (\cite{L13} Proposition 2.4.15).
When $n=4$ the similarity matrix $Z_X(t)$ is invertible for any $t$ $(t>0)$ (\cite{M13} Theorem 3.6 (4)). 
For any finite metric space $X$ the similarity matrix $Z_X(t)$ is invertible for all but finitely many $t>0$ (\cite{L13} Proposition 2.2.6 i). 
Therefore %
$z_X(q)$ $(0<q<1)$ is invertible for any sufficiently small $q$. 


\bigskip
We introduce sets of subscript tuples that play important roles in combinatorial arguments. 
%
\begin{definition}\label{def_I} \rm 
Put 
\[
\mathcal{I}_{k}%
=\displaystyle \left\{(i_0,i_1,\dots,i_k)\,|\,1\le i_0,\dots,i_k \le n, \, 
i_0\ne i_1\ne \dots \ne i_k \right\}  %
\]
for $k\in\bar\NN$, where $\bar\NN=\NN\cup\{0\}$, and 
\begin{equation}\label{eq_def_I}%
\begin{array}{rcl}
\mathcal{I}_\otwop&=&\displaystyle \left\{(i,j,k)\,|\,i\ne j\ne k, \,i<k\right\}, \\[1mm]
%
%

\mathcal{I}_\triangle&=&\{(i,j,k)\,|\,1\le i<j<k\le n\}, \\[2mm]
%
\mathcal{I}_\otp&=& \displaystyle \left\{(i,j,k,l)\,|\,1\le i,j,k,l \le n, \, 
i\ne j\ne k\ne l, \, i<l\right\}, \hspace{1cm}{}  \\[2mm]
%
\mathcal{I}_\sotp&=& \displaystyle \left\{(i,j,k,l)\,|\,1\le i,j,k,l \le n, \, 
i,j,k,l \mbox{ are mutually distinct, } \, i<l\right\}, \hspace{1cm}{}  \\[2mm]
%
\mathcal{I}_{\mbox{\rm \footnotesize disj}}&=&\displaystyle \left\{(i,j,k,l)\,|\,i<j,k<l,\{i,j\}\cap\{k,l\}=\emptyset, \, i<k \right\}, %
%
%
\end{array}
\notag
\end{equation}
where o.$2$-p., o.$3$-p., s.o.$3$-p. and disj stand for open $2$-path, open $3$-path, simple open $3$-path and disjoint respectively. 
When $n=4$, write $\mathcal{I}_{\mbox{\rm \footnotesize disj}}$ as $\mathcal{I}_{\mbox{\rm \footnotesize opp}}$. 
\end{definition}


%
\subsection{Asymptotic behavior at small scale}\label{small-scale_asympt} 
%
Let $\Delta_{X,ij}(t)$ be the $(i,j)$-cofactor of $Z_X(t)$. 
Put 
\[
M\!u_X(t)=\sum_{i,j}\Delta_{X,ij}(t), \quad M\!d_X(t)=\det Z_X(t), 
\] 
then $M_X(t)=M\!u_X(t)/M\!d_X(t)$. 


Let $\nu_k=\nu_k(d_{ij})$ and $\de_k=\de_k(d_{ij})$ be the coefficients of series expansion of $M\!u_X(t)$ and $M\!d_X(t)$ respectively: $\displaystyle M\!u_X(t)=\sum_k\nu_k\,t^k, M\!d_X(t)=\sum_l\de_l\,t^l$. 
Roff and Yoshinaga (\cite{RY} Proof of Theorem 2.3) recently proved\footnote{This can also be verified by direct computation when $n=3,4$.} 
\begin{equation}\label{RY_thm23}
\nu_0=\dots =\nu_{n-2}=0,\> \de_0=\dots =\de_{n-2}=0,\> \nu_{n-1}=\de_{n-1}.
\end{equation}
%
Put $\displaystyle M_\la=\lim_{t\to 0^+}\frac{d^\la}{dt^\la}M_X(t)$ for $\la\ge0$. 
The identities in \eqref{RY_thm23} implies that if $\de_{n-1}\ne0$ then $\displaystyle M_1=\lim_{t\to0^+}M_X'(t)$  is given by 
\begin{equation}\label{M1}
M_1=\frac{\nu_n-\de_n}{\de_{n-1}}.
\end{equation}
%
%
%

The asymptotic behavior of the magnitude function at small scale will be used when $n=3$ and $4$. 
We assume $n=3$ or $4$ in what follows in this Subsection. 
%
\begin{proposition}\label{pos_Mpd}
\begin{enumerate}
\item $\de_2$ is positive for any $3$-point space. 
\item $\de_3$ is non-negative for any $4$-point space and positive if strict inequalities hold in all triangle inequalities. 
\item $\de_3$ is positive for any $4$-point metric subspace of the Euclidean space with the standard metric. 
\end{enumerate}
\end{proposition}

\begin{proof} 
(1) Direct calculation shows 
\begin{equation}\label{cd2}
\de_2=\frac12\sum_{\{i,j,k\}=\{1,2,3\}}(d_{jk}+d_{ik}-d_{ij})(d_{ik}+d_{ij}-d_{jk}).
\end{equation}
%
Among the three terms of the form $d_{jk}+d_{ik}-d_{ij}$ $(i<j)$ at most only one can be $0$. Hence $\de_2>0$. 

\smallskip
(2) Direct calculation shows that $\de_3$ is given by 
\[
-2\sum_{\mathcal{I}_{\mbox{\rm \scriptsize disj}}}\left(d_{ij}^{\,2}d_{kl}+d_{ij}d_{kl}^{\,2}\right)-2\sum_{\mathcal{I}_\triangle}d_{ij}d_{jk}d_{ik}+2\sum_{\mathcal{I}_\scsotp}d_{ij}d_{jk}d_{kl},
\]
and that it is equal to 
\begin{equation}\label{cd3}
\frac16\sum_{\{i,j,k,l\}=\{1,2,3,4\}}(d_{ik}+d_{jk}-d_{ij})(d_{il}+d_{kl}-d_{ik})(d_{ij}+d_{jl}-d_{il}). 
\end{equation}

(3) The above statement implies that $\de_3$ is positive if none of the triangles collapses. 
Suppose that one triangle, say $\triangle P_1P_2P_3$ collapses. 
If $P_4$ is not on the line $L$ through $P_1, P_2$ and $P_3$, then 
\[(d_{24}+d_{12}-d_{14})(d_{34}+d_{23}-d_{24})(d_{14}+d_{13}-d_{34})>0,\]
which implies $\de_3>0$. 
If $P_4$ is on the line $L$, then we may assume without loss of generality that $P_1, P_2, P_3, P_4$ lie on $L$ in this order. Then $\de_3=8d_{12}d_{23}d_{34}>0$. 
\end{proof}
